\section{v0 与 AI-native development：产品约束制造研究问题}

v0 把自然语言想法转换为可运行前端，并逐步扩展到 full-stack application。课堂时点的早期版本使用 GPT-3.5 Turbo，原生 function calling 还不成熟，团队必须构建 bespoke tool invocation。这个历史说明：AI 产品不是简单调用模型 API，而是模型、工具、framework knowledge、preview deployment 和用户反馈的联合系统。

\subsection{从 prompt 到运行中的应用}

前面的 compute-density 讨论关注请求如何高效运行；这里转向 AI application 的生成与验证闭环。本节回答的不是“模型能否写出一段 JSX”，而是一个 prompt 怎样经过 tool invocation、framework-aware code generation、build、preview 和 human feedback 变成可交互产品。读图时应把 preview environment 看成 verifier，而不是单纯的发布终点。

\lecturefigure{14-v0-loop.png}{v0 把生成代码、工具调用、preview URL 和用户反馈连接成产品闭环。}{本地字幕 30:15--31:55；概念重绘。}

\begin{knowledgebox}{读图：运行环境为什么是 coding agent 的 verifier}
代码看起来合理不等于能 build、render 和交互。Preview deployment 提供可执行反馈：编译是否通过、页面是否可访问、视觉和行为是否符合要求。Agent 需要把 build log、browser result 和用户修改重新送回生成循环。
\end{knowledgebox}

\begin{warningbox}{Fast growth 不是 architecture proof}
课堂提到 v0 的快速增长，只能说明需求信号强，不能证明单位经济学、长期 retention 或生成质量已经解决。AI 产品仍需要任务分布、成功率、修复轮数、token cost、latency 和用户留存等长期指标。
\end{warningbox}

\teachervoice{Rauch 把“设定不合理的高标准”当作研究方法：如果目标只是把现有 workflow 自动化，团队可能不会遇到 novel problem；若要求每个 prompt 都快速变成可运行系统，tool calling、sandbox、preview 和 feedback 就会成为新的研究基础设施。}

\subsection{本章小结}

v0 的系统单元不是一个 completion，而是 prompt--tool--code--preview--feedback 循环。产品约束可以创造研究问题，但必须用运行证据验证，而不是只展示生成文本。

